\chapter{Modelagem do problema de perseguição}
\label{cap:modelagem}

\section{Referenciais}
Todo o controle é feito no referencial \code{x500\_depth\_0\_enu}: \textbf{origem no
perseguidor e eixos alinhados ao mundo} (ENU: $x$ leste, $y$ norte, $z$ para cima). Ele
não gira com o drone, então um alvo parado não ``se move'' quando o perseguidor muda o
rumo. A PX4 trabalha em NED; a conversão é $(x,y,z)_{NED} = (y,x,-z)_{ENU}$
(\code{common/frames.hpp}).

\section{Geometria do engajamento}
Sejam $\vect{p}_c, \vect{v}_c$ a posição e a velocidade do perseguidor e
$\vect{p}_t, \vect{v}_t$ as do alvo. Definimos
\begin{equation}
  \rr = \vect{p}_t - \vect{p}_c \quad\text{(posição relativa, \code{p\_rel})},\qquad
  \dot\rr = \vect{v}_t - \vect{v}_c \quad\text{(\code{v\_rel})},\qquad
  \rhat = \rr/\lVert\rr\rVert .
\end{equation}
O objetivo é ficar a uma distância $d$ (\code{target\_distance}, \SI{5}{m}) do alvo,
sobre a linha de visada (LOS). O \emph{ponto de perseguição} (standoff), relativo ao
perseguidor, é (\code{standoffPoint}):
\begin{equation}
  \vect{p}_{des} = \rr - d\,\rhat .
\end{equation}
Quando $\vect{p}_{des}=\vect 0$ o perseguidor está exatamente onde deveria.

\section{Estado de erro}
Os controladores baseados em modelo (LQR, LQI, SMC, MPC) usam o estado de erro
(\code{DoubleIntegrator::errorState}):
\begin{equation}
  \ep = -\vect{p}_{des} \;(\text{perseguidor menos ponto desejado}),\qquad
  \ev = \vect{v}_c - \vect{v}_t .
  \label{eq:erro}
\end{equation}
Derivando $\ep$ (com $d$ constante):
\begin{equation}
  \dot{\vect{e}}_p = -\dot\rr + d\,\dot{\rhat}
  = \ev + \frac{d}{\lVert\rr\rVert}\left(\vect I - \rhat\rhat\T\right)\dot\rr .
  \label{eq:dep}
\end{equation}
O segundo termo é a rotação da LOS. Ele some quando a velocidade relativa é paralela à
LOS. Os modelos lineares o desprezam e usam $\dot{\vect e}_p \approx \ev$. Esse termo
reaparece no \cref{cap:caso}: tratá-lo mal causa uma órbita sem amortecimento.

\section{Modelo duplo integrador}
Se a entrada de controle é a aceleração do perseguidor $\vect u = \vect a_c$ e o alvo tem
aceleração $\vect a_t$ (perturbação),
\begin{equation}
  \dot{\vect e}_p = \ev,\qquad \dot{\vect e}_v = \vect u - \vect a_t .
\end{equation}
Os três eixos são \textbf{desacoplados e idênticos}, então todo o projeto pode ser
feito em \textbf{um eixo} com estado $x = [e_p,\ e_v]\T$:
\begin{equation}
  \dot x = \underbrace{\begin{bmatrix}0&1\\0&0\end{bmatrix}}_{A_c} x
         + \underbrace{\begin{bmatrix}0\\1\end{bmatrix}}_{B_c} u,
  \qquad
  G(s) = \frac{E_p(s)}{U(s)} = \frac{1}{s^2}.
  \label{eq:di}
\end{equation}
Discretizado com segurador de ordem zero (ZOH) e período $\Ts$
(\code{DoubleIntegrator::A,B}):
\begin{equation}
  x_{k+1} = \underbrace{\begin{bmatrix}1&\Ts\\0&1\end{bmatrix}}_{A} x_k
          + \underbrace{\begin{bmatrix}\Ts^2/2\\ \Ts\end{bmatrix}}_{B} u_k,
  \qquad
  G(z) = \frac{\Ts^2}{2}\frac{z+1}{(z-1)^2}.
  \label{eq:did}
\end{equation}

\begin{codigo}
Os controladores que geram aceleração não a enviam ``pura''. O nó integra a aceleração
numa \emph{velocidade de referência} $\vect v_{ref}$
(\code{v\_ref\_ = cmd.velocity}) e publica o par $(\vect v_{ref}, \vect a)$ no
\code{TrajectorySetpoint}. Na prática, a PX4 recebe um setpoint de velocidade com
\emph{feedforward} de aceleração. O estado $\ev$ usado no modelo é
$\vect v_{ref}-\vect v_t$ (\code{state.v\_chaser = v\_ref\_}). Ou seja, o duplo integrador
descreve a \emph{referência}, e o drone real a segue com o atraso da malha interna
(\cref{cap:discreto}).
\end{codigo}

\section{Parâmetros físicos do x500}
Os valores vêm do modelo SDF do PX4 SITL (\code{control/vehicle\_params.hpp}):

\begin{table}[H]
\centering
\caption{Parâmetros do x500 (Gazebo Harmonic) e grandezas derivadas.}
\begin{tabular}{@{}lll@{}}
\toprule
Grandeza & Valor & Origem \\
\midrule
Massa $m$ & \SI{2,064}{kg} & base + 4 rotores \\
Inércia $I_{xx}=I_{yy}$, $I_{zz}$ & \num{0,0217}, \SI{0,04}{kg.m^2} & \code{x500\_base} \\
Constante do motor $k_f$ & \SI{8,549e-6}{N.s^2/rad^2} & \code{MulticopterMotorModel} \\
$\omega_{max}$ / $\omega_{min}$ & 1000 / \SI{150}{rad/s} & \code{SIM\_GZ\_EC\_MAX/MIN} \\
Empuxo máximo $4k_f\omega_{max}^2$ & \SI{34,2}{N} & --- \\
Peso $mg$ & \SI{20,2}{N} & $g=\SI{9,8}{m/s^2}$ \\
Razão empuxo/peso & \num{1,69} & --- \\
Aceleração horizontal máx. (altitude constante) & $g\sqrt{(T/W)^2-1}=\SI{13,4}{m/s^2}$ & --- \\
Constante de tempo do motor & \SI{12,5}{ms} (sobe) / \SI{25}{ms} (desce) & SDF \\
Arrasto linear $D$ & \SI{0,25}{N.s/m} ($m/D\approx\SI{8,3}{s}$) & arrasto dos rotores \\
\bottomrule
\end{tabular}
\end{table}

A aceleração máxima física (\SI{13}{m/s^2}) é muito maior que o limite usado
(\code{max\_acc} = \SI{3}{m/s^2}, inclinação $\approx\arctan(3/9{,}8)=\ang{17}$). Na
malha externa o limite vem da PX4, do conforto da câmera (inclinar demais tira o alvo
do campo de visão) e do atraso, e não dos motores. Com $m/D\approx\SI{8}{s}$, o
arrasto é desprezível na escala de tempo do controle (1--\SI{2}{s}).

\section{Modelo do alvo e do estimador}
O alvo é modelado como de \emph{velocidade constante} com aceleração branca de
intensidade $\sigma_a$ (\code{process\_noise\_std}). O filtro de Kalman
(\code{estimation/constant\_velocity\_kf}) estima $[\rr,\dot\rr]$ a partir da medida de
posição relativa da câmera de profundidade. O nó de controle usa:
\begin{equation}
  \vect v_t \approx \vect v_{c,\text{medida}} + \hat{\dot\rr},\qquad
  \rr(t) \approx \hat\rr(t_k) + \hat{\dot\rr}\,(t-t_k)\ \ \text{(extrapolação pela idade)}.
\end{equation}
\Cref{cap:caso} mostra por que o KF precisa receber a velocidade do próprio perseguidor
como entrada conhecida.
